\documentclass[10pt]{article}
\usepackage[a4paper,margin=1.7cm]{geometry}
\usepackage[T1]{fontenc}
\usepackage{lmodern}
\usepackage{amsmath,amssymb}
\usepackage{longtable,booktabs,array}
\usepackage{listings}
\usepackage[hidelinks]{hyperref}
\lstset{basicstyle=\ttfamily\footnotesize,breaklines=true,columns=fullflexible,keepspaces=true}
\setcounter{MaxMatrixCols}{20}
\allowdisplaybreaks
\setlength{\parindent}{0pt}\setlength{\parskip}{4pt}
\begin{document}

{\LARGE\bfseries dirac16complex: a second-quantized Pin(4,4) fermion field and the 6-planes it creates}\par\medskip

A classroom guide for a beginning student. Everything stated here is computed, with exact arithmetic, in the Mathematica notebook 6planes.nb, which mirrors the notebook Pre-gravity6plane.nb of P. L. Nash. All equations used are written out in this document.


\section*{1. The stage: eight coordinates and a metric with four times}

The coordinates are $x_0, x_1, \dots, x_7$: $x_0$ is the hidden space direction, $x_1, x_2, x_3$ are ordinary 3-space, $x_4$ is time and $x_5, x_6, x_7$ are superluminal time directions. Lengths are measured with


\[
\eta_{AB} = \mathrm{diag}(+1,+1,+1,+1,-1,-1,-1,-1), \qquad ds^2 = \eta_{AB}\,dx^A dx^B = dx_0^2 + dx_1^2 + dx_2^2 + dx_3^2 - dx_4^2 - dx_5^2 - dx_6^2 - dx_7^2 .
\]

\section*{2. The eight real 16x16 Dirac matrices}

With $\sigma = \begin{pmatrix} 0 & I_4\\ I_4 & 0\end{pmatrix}$, the eight real $8\times8$ matrices $\tau_A$ of Pre-gravity6plane.nb ($\tau_0 = I_8$, six antisymmetric matrices built from self-dual and anti-self-dual $4\times4$ blocks, $\tau_7 = \tau_1\cdots\tau_6$) and $\bar\tau_A = \sigma\tau_A^T\sigma$, the Dirac matrices are


\[
T^A \equiv (T16^A)[A] = \begin{pmatrix} 0 & \bar\tau_A \\ \tau_A & 0 \end{pmatrix}, \qquad T^AT^B + T^BT^A = 2\eta^{AB} I_{16} .
\]

All entries are $0$ or $\pm1$ (real). $T^0,\dots,T^3$ are symmetric with square $+I_{16}$; $T^4,\dots,T^7$ are antisymmetric with square $-I_{16}$. Further:


\[
\Gamma = T^0T^1\cdots T^7 = \mathrm{diag}(-I_8, I_8), \qquad \sigma_{16} = T^0T^1T^2T^3 = \mathrm{diag}(-\sigma, \sigma), \qquad S^{AB} = \tfrac14[T^A, T^B],
\]

\[
[S^{AB}, T^C] = \eta^{BC}T^A - \eta^{AC}T^B, \qquad (\sigma_{16}T^A)^T = -\sigma_{16}T^A, \qquad (\sigma_{16}S^{AB})^T = -\sigma_{16}S^{AB}.
\]

The 28 matrices $S^{AB}$ generate Spin(4,4); together with the reflections $T^C$ they generate Pin(4,4).


\section*{3. The field dirac16complex}

\textbf{Definition.} dirac16complex is a 16-component Pin(4,4) spinor field


\[
\psi(x) = (\psi_0(x), \psi_1(x), \dots, \psi_{15}(x))^T, \qquad \psi_a = \tfrac{1}{\sqrt2}(\xi_a + i\chi_a),
\]

whose components are complex and anticommuting (Grassmann-odd): with $\xi_a, \chi_a$ real Grassmann generators,


\[
\psi_a\psi_b = -\psi_b\psi_a, \qquad \psi_a\psi_b^{*} = -\psi_b^{*}\psi_a, \qquad \psi_a^2 = 0, \qquad (\theta\eta)^{*} = \eta^{*}\theta^{*}.
\]

The last rule (complex conjugation reverses the order of factors) is the classical shadow of Hermitian conjugation of operators. dirac16complex is a second-quantized field: the $\psi_a$ are field operators, and the Grassmann variables are their classical (path-integral) counterparts. A commuting (c-number) solution of the field equation is a mode function of the field; following Pre-gravity6plane.nb we call such a mode function the wave function of a Universe.


\textbf{Dirac adjoint.} $\bar\psi = \psi^\dagger\sigma_{16}$. Under Spin(4,4), $\psi \to g\psi$ with $g = \exp(\tfrac12\omega_{AB}S^{AB})$, the scalar $\bar\psi\psi$ is invariant and $\bar\psi T^A\psi$ is a vector.


\section*{4. The Lagrangian: kinetic, mass (quadratic) and self-coupling (quartic) terms}

\[
\mathcal L = \tfrac12\left(\bar\psi\,T^A\partial_A\psi - (\partial_A\bar\psi)\,T^A\psi\right) - m\,\bar\psi\psi - \frac{\lambda}{2}\,(\bar\psi\psi)^2 ,
\]

with $m \neq 0$ and $\lambda \neq 0$ real constants and $A$ summed over $0,\dots,7$. Compare with the Higgs potential $V = \mu^2\lvert\phi\rvert^2 + \lambda\lvert\phi\rvert^4$: here the mass term is quadratic and the self-coupling quartic in the spinor components; the field is a fermion, so $(\bar\psi\psi)^2$ is the simplest non-trivial self-interaction. Proven properties (exact Grassmann algebra, 392 monomials):


\begin{itemize}
\item $\mathcal L$ is real: $\mathcal L^{*} = \mathcal L$ exactly (no factor $i$ is needed because $\sigma_{16}T^A$ is real antisymmetric).
\item The mass and quartic terms are non-zero because the components are complex. For a real (Majorana) field $\xi^T\sigma_{16}\xi = 0$ identically, since $\sigma_{16}$ is symmetric and $\xi_a\xi_b = -\xi_b\xi_a$.
\item $\mathcal L$ is Spin(4,4) invariant, and invariant under the phase rotation $\psi \to e^{i\alpha}\psi$.
\end{itemize}

\section*{5. The Euler-Lagrange equations (flat space)}

Treat $\psi_a$ and $\psi_a^{*}$ as independent Grassmann variables and vary $\psi_a^{*}$ with left derivatives:


\[
\frac{\partial^L\mathcal L}{\partial\psi_a^{*}} - \partial_\mu\,\frac{\partial^L\mathcal L}{\partial(\partial_\mu\psi_a^{*})} = 0 \quad (a = 0,\dots,15).
\]

The exact computation gives the 16 equations


\[
\sigma_{16}\left(T^A\partial_A\psi - m\psi - \lambda(\bar\psi\psi)\psi\right) = 0 \quad\Longleftrightarrow\quad T^A\partial_A\psi = (m + \lambda\,\bar\psi\psi)\,\psi ,
\]

and varying $\psi_a$ gives the complex-conjugate equation $(\partial_A\bar\psi)T^A = -(m + \lambda\bar\psi\psi)\bar\psi$. Example (first component, $\lambda = 0$):


\[
m\psi_4 - \partial_0\psi_{12} + \partial_1\psi_{11} - \partial_2\psi_{10} + \partial_3\psi_9 - \partial_4\psi_9 + \partial_5\psi_{10} + \partial_6\psi_{11} + \partial_7\psi_{12} = 0 .
\]

\textbf{Conserved current.} $J^A = \bar\psi T^A\psi$ satisfies $\partial_AJ^A = \bar\psi(T^A\partial_A\psi) + (\partial_A\bar\psi)T^A\psi = (m+\lambda\bar\psi\psi)\bar\psi\psi - (m+\lambda\bar\psi\psi)\bar\psi\psi = 0$.


\section*{6. Second quantization and the indefinite (Krein) metric}

Quantize with $x_4$ as time. The momentum conjugate to $\psi$ is $\pi = \bar\psi T^4$, and $\{\psi_a(x), \pi_b(y)\} = i\delta_{ab}\delta(x-y)$ gives


\[
\{\psi_a(x), \psi_b^\dagger(y)\} = K_{ab}\,\delta(x-y), \qquad K = -i\,\sigma_{16}T^4, \qquad K^\dagger = K, \qquad K^2 = I_{16},
\]

with eight eigenvalues $+1$ and eight $-1$. The space of states therefore carries an indefinite (Krein) metric. This cannot be avoided: $g_{45} = T^4T^5$ (the rotation by $\pi$ in the $(x_4,x_5)$-plane, an element of Spin(4,4)) obeys $g_{45}^T(\sigma_{16}T^4)g_{45} = -\sigma_{16}T^4$, so it reverses the sign of the $x_4$-charge density $\bar\psi T^4\psi$; with four time directions there is no invariant notion of positive energy.


\textbf{Two mirror sectors.} Since $\sigma_{16} = K\beta$ with $\beta = -iT^4$, and $K$ commutes with $T^0, \dots, T^4$ (and anticommutes with $T^5, T^6, T^7$), every field that does not depend on $x_5, x_6, x_7$ splits as $\psi = \psi_+ + \psi_-$, $K\psi_\pm = \pm\psi_\pm$, and


\[
\mathcal L = \mathcal L_D[\psi_+] - \mathcal L_D[\psi_-] - \frac\lambda2\,(S_+ - S_-)^2, \qquad \mathcal L_D[\chi] = \chi^\dagger\beta(T^A\partial_A - m)\chi\ (A=0..4), \qquad S_\pm = \psi_\pm^\dagger\beta\psi_\pm .
\]

The 16-component field is thus a pair of 8-component universes whose energies (masses) enter with opposite signs: a "+mass" and a "-mass" universe, coupled by $-\lambda S_+S_-$ inside $-\frac\lambda2(S_+-S_-)^2 = -\frac\lambda2(S_+^2 + S_-^2) + \lambda S_+S_-$.


\section*{7. Mirroring Pre-gravity6plane.nb: dirac16complex depending on $x_0$ and $x_4$}

Pre-gravity6plane.nb studies $\Psi_{16} = f16[k][x_0, x_4]$, a wave function of a Universe depending only on $x_0$ and $x_4$, with Lagrangian $L_a = \frac1H\Psi_{16}^T\sigma_{16}\sum_A T^A\partial_A\Psi_{16} + \frac{2M}{H}\Psi_{16}^T\sigma_{16}\Psi_{16}$ and field equation $T^A\partial_A\Psi_{16} + 2M\Psi_{16} = 0$. The corresponding configurations of dirac16complex are those that depend only on $x_0$ and $x_4$ (no dependence on $x_1, x_2, x_3, x_5, x_6, x_7$). For them the field equation is


\[
T^0\partial_0\psi + T^4\partial_4\psi = (m + \lambda\bar\psi\psi)\psi, \qquad \text{with } m = -2M \text{ to match the old notebook.}
\]

\textbf{Couplings (analogue of eLaCouplings).} The 16 equations split into four independent systems of four components:


\[
\{0,5,8,13\},\quad \{1,4,9,12\},\quad \{2,7,10,15\},\quad \{3,6,11,14\},
\]

exactly the pattern found in Pre-gravity6plane.nb.


**Variables $z = 6Hx_0$, $t = Hx_4$ (analogue of eLazt).** The linear equation becomes $6T^0\partial_z\psi + T^4\partial_t\psi = \frac{m}{H}\psi$. A mode $e^{pz+qt}u$ solves it iff $(6pT^0 + qT^4 - \frac mH)u = 0$, possible iff


\[
(6p)^2 - q^2 = \left(\frac mH\right)^2 \quad\Longleftrightarrow\quad (6p+q)(6p-q) = \frac{m^2}{H^2}.
\]

Writing $6p + q = 2C$ and $m = -2M$, the exponent is


\[
pz + qt = C\left(t + \frac z6\right) - \frac{M^2}{CH^2}\left(t - \frac z6\right),
\]

which is exactly the family of exponential solutions that Maple produced in Pre-gravity6plane.nb ("Solution 15feb"); the field equation then reads $(6pT^0 + qT^4 + \frac{2M}{H})u = 0$ and the spinor is $u = (6pT^0 + qT^4 - \frac{2M}{H})w$ for any constant $w$, because $(6pT^0+qT^4+\frac{2M}H)(6pT^0+qT^4-\frac{2M}H) = (36p^2-q^2-\frac{4M^2}{H^2})I_{16} = 0$. For a complex field the generic solutions are oscillatory plane waves:


\[
\psi = e^{i(kx_0 - Ex_4)}\left(ikT^0 - iET^4 + m\right)w, \qquad E = \pm\sqrt{k^2+m^2}.
\]

\textbf{Retarded and advanced variables.} $z_{adv} = 6t + z$, $z_{ret} = 6t - z$ (equivalently $\xi_{adv} = x_0 - x_4$, $\xi_{ret} = x_0 + x_4$).


\textbf{Light-like quanta.} For $m = 0$ the solution $\psi = e^{ik(x_0-x_4)}\,ik(T^0 - T^4)w$ depends on one light-cone variable only. Because $(T^0 - T^4)^2 = (T^0)^2 + (T^4)^2 = 0$, every component of the current $J^A = i\bar\psi T^A\psi$ and the scalar $\bar\psi\psi$ vanish identically, just as the bilinears $\psi\sigma_{16}\psi$ vanished for the light-like quanta of Pre-gravity6plane.nb.


\section*{8. Triality frames and vector forms}

The unit spinor is $u = \frac1{\sqrt2}(1,0,0,0,1,0,0,0)$, the 8th normalized eigenvector of $\sigma$ ($u^T\sigma u = 1$). The triality frames of Pre-gravity6plane.nb are


\[
(F_A{}^a)_{aA} = \eta_{AA}(\tau_A u)_a, \qquad (F_a{}^A)_{Aa} = (u^T\sigma\bar\tau_A)_a, \qquad F_A{}^aF_a{}^A = F_a{}^AF_A{}^a = I_8 .
\]

For a Universe $\psi$ write $\psi_1 = (\psi_0,\dots,\psi_7)$ (type 1) and $\psi_2 = (\psi_8,\dots,\psi_{15})$ (type 2), and define the vector forms


\[
V_2 = F_a{}^A\,\psi_2, \qquad V_1 = \eta\,(\psi_1^T\sigma F_A{}^a) .
\]

\section*{9. The 6-planes of dirac16complex}

The field equation is real, so the real part of every Universe that depends only on $x_0$ and $x_4$ is itself a real solution of the type studied in Pre-gravity6plane.nb. Apply the old construction to it: $n_1 = V_1[\mathrm{Re}\,\psi]$ and $n_2 = V_2[\mathrm{Re}\,\psi]$ are real 8-vectors depending on $(x_0, x_4)$, and with $Z = (z_1,\dots,z_8)\in\mathbb R^8$:


\[
\text{7-plane}_k = \{Z : n_k\cdot Z = 0\}, \qquad M^6 = \text{7-plane}_1 \cap \text{7-plane}_2 = \{Z : n_1\cdot Z = 0,\ n_2\cdot Z = 0\}.
\]

$M^6$ is 6-dimensional whenever $n_1$ and $n_2$ are linearly independent. Solving for a pair of coordinates $\{z_i, z_j\}$ is possible iff the minor $n_{1i}n_{2j} - n_{1j}n_{2i} \neq 0$; there are $\binom82 = 28$ such pairs, the 28 parameterizations of $M^6$. \textbf{Example} (computed in 6planes.nb): for $\psi = e^{i(\frac34x_0 - \frac54x_4)}(\frac34 iT^0 - \frac54 iT^4 + 1)w$ with $w = (1,2,\dots,16)$ at $(x_0,x_4) = (0.3, 0.7)$, the rank of $\{n_1, n_2\}$ is 2 and all 28 minors are non-zero, and


\[
z_7 = -132.890\,z_1 + 189.925\,z_2 - 174.913\,z_3 + 155.902\,z_4 + 77.0462\,z_5 - 74.0462\,z_6,
\]

\[
z_8 = -149.525\,z_1 + 213.689\,z_2 - 196.301\,z_3 + 174.913\,z_4 + 87.5520\,z_5 - 84.5520\,z_6 .
\]

\textbf{A pair of universes creates a 6-plane by itself.} For a single plane wave, $V_2 = e^{i\phi}\times$(real vector), so $\mathrm{Re}\,V_2$ and $\mathrm{Im}\,V_2$ are parallel and only give a 7-plane. For the superposition of a positive-energy and a negative-energy Universe,


\[
\psi_{pair} = e^{i(kx_0 - Ex_4)}(ikT^0 - iET^4 + m)w + e^{i(kx_0 + Ex_4)}(ikT^0 + iET^4 + m)w', \qquad E = \sqrt{k^2+m^2},
\]

$\mathrm{Re}\,V_2$ and $\mathrm{Im}\,V_2$ are independent (rank 2) and


\[
M^6_{pair} = \{Z : \mathrm{Re}\,V_2\cdot Z = 0,\ \mathrm{Im}\,V_2\cdot Z = 0\}
\]

is a 6-plane that moves with $(x_0, x_4)$: the pair $\{+E, -E\}$ of universes generates a 6-dimensional plane.


\section*{10. What a student should remember}

\begin{enumerate}
\item The Dirac matrices are real because the signature is (4,4); Cl(4,4) is the algebra of all real $16\times16$ matrices.
\item The field must be complex for $\bar\psi\psi$ and $(\bar\psi\psi)^2$ to be non-zero.
\item The field equation $T^A\partial_A\psi = (m+\lambda\bar\psi\psi)\psi$ follows from the Grassmann Euler-Lagrange equations.
\item With four times, quantization has an indefinite metric and the field splits into two mirror universes with opposite energies.
\item Universes (mode functions) produce 7-planes through their vector forms; two 7-planes, or a single $\pm$ pair of universes, produce a 6-plane $M^6$.
\end{enumerate}

\end{document}